\documentclass[exercice,a4]{thedocument}%
\usepackage{texmaths-path}%
\usepackage{texmaths-tikz}%
\usepackage{mdframed}%
\usepackage{enumitem}%
\startdatakeys
\source{19GENMATAMS}
\cycle{4}
\competencesDuSocle{CA,RA,CO,CH}
\connaissancesRequises{D2e1,D2f*}
\competencesTravaillees{A1i*,A1g*,D2d*,D2e*}
\enddatakeys
\begin{document}%
\!scale=1.4;%
Un silo à grains permet de stocker des céréales. Un ascenseur permet d'acheminer
le blé dans le silo.\par\noindent L'ascenseur est soutenu par un pilier.%
\par\noindent
{\centering\pic[scale=0.3]{19GENMATAMS-ex4-fig1.png}\par}\medskip
\begin{minipage}[t]{.5\linewidth}
    On modélise l'installation par la figure ci-dessous qui n'est pas réalisée
    à l'échelle~:
    \begin{enumerate}
        \item Les points \mathspoint{C}, \mathspoint{E} et \mathspoint{M}
        sont alignés;
        \item Les points \mathspoint{C}, \mathspoint{F}, \mathspoint{H}
        et \mathspoint{P} sont alignés;
        \item Les droites \mathspoint{(EF)} et \mathspoint{(MH)} sont
        perpendiculaires à la droite \mathspoint{(CH)};
        \item $\mathspoint{CH}=8,50$ m et $\mathspoint{CF}=2,50$ m;
        \item Hauteur du cylindre : $\mathspoint{HM}=20,40$ m;
        \item Diamètre du cylindre : $\mathspoint{HP}=4,20$ m. 
    \end{enumerate}
\end{minipage}
\begin{minipage}[t]{.4\linewidth}
    \begin{tikzpicture}[scale=\*scale;,baseline=(current bounding box.north)]
        \coordinate[label=below:C](C)at(0,0);
        \coordinate[label=below:F](F)at(1,0);
        \coordinate[label=below:H](H)at(3,0);
        \coordinate[label=below:P](P)at(4,0);
        \coordinate[label=above left:E](E)at(1,1);
        \coordinate[label=above:M](M)at(3,3);
        \coordinate(Q)at(4,3);
        \draw[fill=\currentcolor](H)--(P)--(Q)--(M)--cycle;
        \tkzMarkRightAngle[colored](H,F,E)
        \tkzMarkRightAngle[colored](P,H,M)
        \tkzDrawLine[add=0.2 and 0.2](C,P)
        \draw(C)--(M);\draw(E)--(F);
        \path(M)--(P)node[midway]{\footnotesize Cylindre};
    \end{tikzpicture}
\end{minipage}\par\bigskip
\noindent{\textbf{Les quatres questions suivantes sont indépendantes.}}
\begin{enumerate}
    \item Quelle est la longueur \mathspoint{CM} de l'ascenseur à blé~?
    \item Quelle est la hauteur \mathspoint{EF} du pilier~?
    \item Quelle est la mesure de l'angle $\widehat{\mathspoint{HCM}}$
    entre le sol et l'ascenseur à blé~?\par\noindent On donnera une
    valeur approchée au degré près.
    \item Un mètre cube de blé pèse environ 800 kg.\par\noindent
    Quelle masse maximale de blé peut-on stocker dans ce silo~? On
    donnera la réponse à une tonne près.
\end{enumerate}
{\centering
\begin{minipage}{.6\linewidth}
    \begin{mdframed}
        \noindent\textbf{Rappels :}
        \begin{itemize}[label=--]
            \item 1 tonne = 1\,000 kg;
            \item volume d'un cylindre de rayon \mathspoint{R} et de hauteur
            \mathspoint{h} : $\pi\times\mathspoint{R}^2\times\mathspoint{h}$;
        \end{itemize}
    \end{mdframed}
\end{minipage}
\par}
\showthesource
\end{document}